\chapter{Physical Detection Systems}
\label{ch:2.4}
\chaptertopic{Topic 2.4}

\chapterquote{Trust, but verify.}{Russian proverb, popularized in English by Ronald Reagan}

\section{Motivation}

Chapter~\ref{ch:2.3} was about controls that keep an attacker out. This closing chapter of Unit~2
is about the layer that assumes a control has already failed, or is being actively tested, and
asks a narrower question: how does a facility \emph{notice}? Access logs, alarm systems, and
motion sensors are all detective, not preventive, controls -- none of them stops an intruder
physically. Their value is entirely in how fast and how reliably they shrink the gap between an
intrusion starting and someone finding out. That gap is exactly how much damage an
attacker gets to do uncontested.

\begin{apalignment}{Unit 2: Securing Spaces}{Topic 2.4 -- Detecting Physical Attacks}
\begin{apobjectives}
\apobjective{2.4.A}{Detect Attacks}{Identify the ways a security control can reveal a physical attack -- the sensor that reports a state change, the camera that supplies the assessment, the log that records a credential, and the human patrol that observes what none of them are positioned to see.}
\apobjective{2.4.B}{Detect Attacks}{Determine where a detective control has to sit to be useful: which boundary, which choke point, which asset, and why a sensor placed at the wrong point in an intruder's path detects the attack too late to act on it.}
\apobjective{2.4.C}{Detect Attacks}{Apply detection techniques to a specific situation, reading the evidence a physical intrusion leaves across several independent systems and assembling it into an account that will survive scrutiny.}
\end{apobjectives}
\end{apalignment}

\section{Access Logs and Monitoring}

Every badge reader, key-card door, and biometric scanner generates a continuous stream of
timestamped records the instant someone attempts or succeeds at entry.

\begin{definition}{Physical access log}
A timestamped record of credential presentations at controlled entry points, typically including
who presented a credential, which door, when, and whether it was granted or denied.
\end{definition}

Unlike a camera feed, which requires a human to watch, a log can be searched and analyzed
automatically -- which is what makes it a genuinely active detective control rather than just a
record. Three recurring patterns distinguish an attack from ordinary use. A \textbf{location
inconsistency} is a single credential appearing at two geographically separated readers within a
window too short for legitimate travel. The same badge logged at a front lobby at 07:55 and at a
server-room door at 02:17 the same night is not one person moving quickly -- it is a cloned or
shared credential. A \textbf{repeated denial burst} is a cluster of ACCESS DENIED events for one
credential in a short window. This is the access-log equivalent of a network brute-force attack, and
detected with the same logic: count failures per credential per window, alert past a threshold.
\textbf{After-hours access} requires a baseline defined per door, not per building -- a server
room is expected to be empty at midnight even in a building with round-the-clock cleaning staff.
Any access to it outside its authorized schedule warrants an alert, regardless of whether the
badge itself is valid.

\begin{lstlisting}[language=Python, caption={Flagging access-log anomalies across three rules}]
def flag_access_anomalies(entries, denial_threshold=5, denial_window=60,
                           location_gap_seconds=300, after_hours_start=22,
                           after_hours_end=6, restricted_doors=frozenset()):
    """entries: list of dicts like {"badge_id": str, "door": str,
       "timestamp": int, "hour": int, "result": str}.
       Returns sorted badge IDs triggering any alert rule."""
    flagged = set()
    by_badge = {}
    for e in entries:
        by_badge.setdefault(e["badge_id"], []).append(e)

    for badge_id, badge_entries in by_badge.items():
        badge_entries.sort(key=lambda e: e["timestamp"])
        denials = [e["timestamp"] for e in badge_entries if e["result"] == "DENIED"]
        left = 0
        for right in range(len(denials)):
            while denials[right] - denials[left] > denial_window:
                left += 1
            if right - left + 1 >= denial_threshold:
                flagged.add(badge_id)

        for e in badge_entries:
            if e["door"] in restricted_doors:
                if e["hour"] >= after_hours_start or e["hour"] < after_hours_end:
                    flagged.add(badge_id)

        for a, b in zip(badge_entries, badge_entries[1:]):
            if a["door"] != b["door"] and b["timestamp"] - a["timestamp"] < location_gap_seconds:
                flagged.add(badge_id)
    return sorted(flagged)
\end{lstlisting}

Tailgating (Chapter~\ref{ch:2.2}) is the one pattern a badge log cannot catch directly, because the
attacker never swipes at all. It can still surface indirectly, though: if the count recorded by an
anti-passback badge system falls below what an independent people-counter sensor observed at the
same door, the difference itself is the alert. In practice, an organization aggregates all of
this into a Security Information and Event Management (SIEM) platform that applies threshold,
schedule, and correlation rules across every reader at once. The platform then ranks the resulting alerts by
severity for a tier-1 analyst to triage.

\begin{misconception}
\textbf{``A valid badge swipe means the entry was authorized.''} A badge can be cloned, stolen, or
knowingly shared; a valid swipe only proves a credential was presented, not that the right person
presented it. Context -- the time, the location, the frequency -- is what actually determines
legitimacy.
\end{misconception}

\section{Alarm Systems}

An alarm system occupies the middle ground between a lock that keeps someone out and a camera
that merely records that they got in anyway: it is a detective control whose job is to notice an
intrusion in progress and trigger a response fast enough to matter. Nothing about a siren
physically stops someone from climbing through a broken window. Treating an armed alarm as
equivalent to a locked door is the single most common physical-security planning mistake, because
it misallocates budget and attention toward a control that cannot actually prevent anything.

An alarm system has four parts. \textbf{Sensors} detect a specific event at a specific
location: a door/window contact (a reed switch that opens when the door swings), or a
glass-break sensor tuned specifically to the acoustic signature of shattering glass, to avoid
false trips from ordinary loud noise. A \textbf{panic button} is different in kind: it is not
triggered by an intrusion event at all, but by a person deliberately signaling they need help. The
\textbf{control panel} tracks which zones are armed, runs timing logic, and decides whether a trip
sounds a local siren, calls the monitoring center, or both. The \textbf{monitoring / central
station} is a staffed facility that receives the panel's signal, attempts to verify it, and
dispatches police if warranted.

\begin{figure}[h]
\centering
\begin{tikzpicture}[
  box/.style={draw=kickstartermaroon, thick, rounded corners, minimum width=2.6cm, minimum height=0.9cm, align=center, font=\scriptsize},
]
\node[box, fill=kickstartercream] (sensor) at (0,0) {Sensor trips};
\node[box] (panel) at (3,0) {Control panel:\\zone armed? in delay?};
\node[box, fill=kickstartermaroon!10] (susp) at (6.4,1.3) {Suppressed\\(disarmed / in delay)};
\node[box, fill=kickstartermaroon!10] (local) at (6.4,0) {Local-only\\(siren, no dispatch)};
\node[box, fill=kickstartermaroon!25] (dispatch) at (6.4,-1.3) {Verified dispatch\\(cross-zone / panic)};
\draw[-{Latex[length=2mm]}] (sensor) -- (panel);
\draw[-{Latex[length=2mm]}] (panel) -- (susp);
\draw[-{Latex[length=2mm]}] (panel) -- (local);
\draw[-{Latex[length=2mm]}] (panel) -- (dispatch);
\end{tikzpicture}
\caption{Response is tiered, not binary: the same sensor trip resolves to suppression,
a local-only siren, or a verified dispatch depending on arm state, delay windows, and
corroborating signals.}
\label{fig:2.4-alarm-tiers}
\end{figure}

\begin{definition}{Entry/exit delay}
A timed grace period -- typically 45 to 60 seconds -- that lets a legitimate occupant leave through
a designated door after arming, or reach the panel to disarm on arrival, without a sensor trip
during that window counting as an intrusion.
\end{definition}

The hardest engineering problem in alarm design is not detecting a sensor trip; it is telling a
trip caused by an intruder apart from one caused by the occupant who just armed the system and is
walking out the door. That distinction is exactly what entry/exit delay windows solve. False alarms are the alarm
industry's central operational problem: the overwhelming majority of dispatches trace back to
pets, wind-blown curtains, or a door not fully latched. Many jurisdictions now fine or
decline to dispatch to repeat false-alarm addresses as a result. \textbf{Cross-zone verification} -- requiring
two independent sensors to trip within a short window of each other before treating an event as
high-confidence -- and \textbf{video verification} -- letting a monitoring-center operator visually
confirm a person before calling police -- both reduce false dispatch. Systems supporting
verification typically receive faster, higher-priority response, because dispatchers can trust the
alert.

\begin{lstlisting}[language=Python, caption={Classifying alarm events by confidence and response tier}]
def classify_alarm_event(sensor_type, arm_state, in_exit_delay, in_entry_delay,
                          cross_zone_confirmed):
    """Returns one of "suppressed", "false_alarm", "local_only", "dispatch"."""
    if sensor_type == "panic":
        return "dispatch"
    if arm_state == "disarmed":
        return "suppressed"
    if in_exit_delay or in_entry_delay:
        return "suppressed"
    if cross_zone_confirmed:
        return "dispatch"
    return "local_only"
\end{lstlisting}

Response is deliberately tiered. A local-only response -- siren and strobe, no call to
authorities -- suits low-confidence, single-sensor trips and deters an opportunist without
consuming monitoring-center resources on every ambiguous trip. A verified dispatch is reserved for
high-confidence events. Panic and duress signals bypass every other layer of this logic entirely
-- no verification, no arm-state check -- and escalate straight to dispatch, because a human
deliberately pressed the button. The cost of ignoring a real emergency far outweighs the cost
of occasionally dispatching on a false one.

\begin{misconception}
\textbf{``An armed alarm prevents a break-in.''} An alarm is a detective and response control, not
a preventive one -- it can only notice and react to an intrusion already underway, not physically
stop it. That job belongs to the locks, mantraps, and barriers of Chapter~\ref{ch:2.3}.
\end{misconception}

\begin{example}[Why an excellent detector still cries wolf almost every time]
Consider a zone with a genuinely good detector: it alarms on a real intrusion 99 times out of 100.
On any ordinary night with no intruder, it produces a nuisance alarm 1 percent of the time.
Suppose a real intrusion at this zone happens on roughly one night in a thousand. Across a
thousand nights, the detector fires about one true alarm from the single real event and about ten
false alarms from the 999 quiet nights (0.01~$\times$~999~$\approx$~10). An alarm from this zone is
therefore genuine roughly $1/(1+10) \approx 9\%$ of the time -- despite every published
specification of the sensor being excellent. Nothing is wrong with the hardware; the event's
rarity, not the detector's quality, dominates the result. This is the same base-rate reasoning that
governs any low-prevalence screening problem, from medical testing to network intrusion detection.
It also has a real historical price tag. Problem-oriented policing research places the false rate
for police-dispatched burglar alarms at roughly 94 to 98 percent. In 1999, only about 0.3 percent
of the alarm calls Salt Lake City police attended turned out to involve an actual
crime.\srccite{2.4}{1} In 2000, Salt Lake City adopted a verified-response
policy under which police attend an alarm only after a private responder or camera has confirmed a
crime in progress, and alarm-related dispatch volume fell sharply as a result.
\end{example}

Three practical conclusions follow directly from the arithmetic above. First, a monitoring
operation cannot be designed on the assumption that alarms are usually real, because for rare
events they structurally never will be. It has to be designed instead to dispose of the common (false)
case cheaply. Second, the way to raise the fraction of genuine alarms is to add an
\emph{independent} second signal, not to make one sensor stricter. This is exactly why cross-zone
verification and the dual-technology AND-logic covered later in this chapter work: they attack the
ten false alarms rather than trying to make the one real alarm more certain. Third, the cost of
false alarms is not a minor operational nuisance. It is the mechanism by which a detection program
stops being answered at all, whether the responder is a police dispatcher or a
security-operations-center analyst deciding which of a thousand daily alerts deserves a look.

\section{Motion Sensors}

A badge reader only fires when a badge is presented, and no organization watches every camera feed
in real time. Motion sensors close exactly that gap: always-on detective controls that alert the
instant something moves in a space that should be empty, without identifying who moved or why.

\begin{definition}{Passive infrared (PIR) sensor}
A motion sensor that emits no energy of its own and instead watches the infrared (heat) radiation
already present in a room. It divides its field of view into segments and detects a rapid
\emph{change} in infrared level as a warm object crosses a segment boundary.
\end{definition}

Because PIR responds to \emph{change}, not to heat itself, a person standing perfectly still long
enough can occasionally go undetected -- no motion means no delta, and no delta means no trigger.
\textbf{Microwave sensors} work differently: they actively emit a signal and measure its
reflection, detecting motion via the \textbf{Doppler shift} in the returning wave -- the same
physics behind a radar speed gun. Microwave energy passes through thin walls and glass that block
infrared. That is a strength for coverage and a weakness for containment, since a microwave sensor can
trigger on motion in an adjacent room outside the zone it is meant to protect. Ultrasonic sensors
apply the same Doppler principle with sound instead, at shorter range and with more susceptibility
to HVAC-airflow false triggers.

\begin{example}[Why dual-technology sensors use AND, not OR]
A PIR element triggers on sudden heat-source movement that is not a person -- sunlight sweeping a
wall, a heating vent kicking on. A microwave element triggers on motion a PIR would never see --
vibration from a passing truck, HVAC airflow. Combining both technologies in one housing and
requiring \emph{both} to trigger within the same short window, rather than either alone, sharply
cuts false positives. Sunlight glare fools a PIR element but produces no Doppler shift, and
airflow can fool a microwave element but produces no meaningful infrared delta. The tradeoff is a
small risk of missing a true detection, if an intruder's movement happens not to trigger both
elements simultaneously. That is an acceptable cost in a zone where every false dispatch has a real
price.
\end{example}

\begin{lstlisting}[language=Python, caption={Classifying a motion-sensor trigger as a true detection or false positive}]
KNOWN_FALSE_CAUSES = {
    "pir": {"sunlight", "hvac_vent", "pet", "insect", "curtain_sway"},
    "microwave": {"hvac_airflow", "adjacent_room_traffic", "radio_interference"},
}

def classify_trigger(sensor_type, cause, dual_tech_confirmed=None):
    """sensor_type: "pir", "microwave", or "dual". cause: str describing
       what triggered the sensor. dual_tech_confirmed: bool, only relevant
       when sensor_type == "dual" -- whether both elements agreed."""
    if sensor_type == "dual":
        return "true_detection" if dual_tech_confirmed else "false_positive"
    if cause in KNOWN_FALSE_CAUSES.get(sensor_type, set()):
        return "false_positive"
    return "true_detection"
\end{lstlisting}

Placement mirrors camera placement from Chapter~\ref{ch:2.3}: mount high enough to resist
tampering, and angle the detection cone toward doors and choke points rather than open floor space.
Avoid pointing a PIR sensor at a window, an HVAC vent, or anything that sways in an air
current, since each is a documented false-trigger source. \textbf{Sensitivity tuning} is a
tradeoff dial, not a one-directional improvement. Higher sensitivity catches slower, more careful
movement, but also a mouse or a shifting shadow -- and enough nuisance triggers train staff to ignore
alerts entirely, which is a worse outcome than a slightly less sensitive sensor. Facilities
compound this with \textbf{zoned detection}: separate sensors covering distinct areas so an
alert also identifies \emph{where}. Perimeter sensors (fences, exterior contacts) layer
outside interior sensors (room-level PIR and dual-technology units in high-value rooms), giving
defense in depth if the perimeter layer is defeated.

\begin{misconception}
\textbf{``Higher sensitivity is always better security.''} Maximum sensitivity also catches pets,
insects, and sunlight glare as false positives, and enough of those get a sensor's alerts
deprioritized or disabled by fatigued staff -- the opposite of better security.
\end{misconception}

\section*{Review Questions}
\begin{reviewquestions}
\item A badge shows a valid entry at the front lobby at 07:55 and a valid entry at a server-room
door 250 miles away at 08:03 the same morning. Explain which detection rule this triggers and why
the badge being ``valid'' at both doors does not make the access legitimate.
\item Explain why an alarm system's entry/exit delay is necessary even though its purpose,
described plainly, is to \emph{not} alert on some real sensor trips.
\item A monitoring center receives a single door-contact trip with no corroborating sensor. Explain
why cross-zone verification would change how this event is classified, and what specific risk
verification is designed to reduce.
\item Explain why a dual-technology motion sensor uses AND-logic between its PIR and microwave
elements rather than OR-logic, and describe the one scenario in which this design choice could
cause a missed true detection.
\item A facility experiences frequent nuisance alerts from a PIR sensor mounted near a window.
Using the placement and sensitivity concepts from this chapter, explain two independent changes
that would reduce false positives without simply disabling the sensor.
\item In the worked base-rate example, a detector that is 99\% accurate on both true and false
alarms is still only genuine about 9\% of the time it fires. Explain, without recomputing the
arithmetic, why adding a second, independent sensor and requiring both to agree raises this
fraction far more effectively than replacing the single sensor with an even more accurate one.
\item Salt Lake City's verified-response policy means police attend an alarm only after a private
responder or camera confirms a crime in progress. Explain what problem in the alarm-response
pipeline this policy directly targets, and why it is a structural fix rather than merely a stricter
threshold on the same detector.
\end{reviewquestions}

\begin{answerkey}
\item This triggers a location-inconsistency rule: the same credential appearing at two
geographically separated readers within a window too short for legitimate travel. A badge being
``valid'' at both doors only proves a credential was presented and accepted at each reader
individually -- it says nothing about whether the same physical person carried it to both places.
250 miles in eight minutes is physically impossible for one person, which is exactly the
context (time and location together) that the chapter argues a raw badge-valid signal cannot
supply on its own.
\answerlabel{2.4.A}
\item The entry/exit delay is necessary precisely because a legitimate occupant arming the system
and walking out, or arriving and walking to the panel to disarm, produces a sensor trip that is
indistinguishable at the sensor level from an intruder's trip. Without the delay window, the system
cannot tell the two apart, and every legitimate arming or disarming cycle would itself trigger a
false alarm. Deliberately not alerting on that specific, predictable, low-risk trip is what makes
the rest of the system's alerts trustworthy.
\answerlabel{2.4.B}
\item A single door-contact trip with no corroborating sensor is low-confidence: it could be an
intrusion, or a door not fully latched, wind, or routine staff movement outside the expected
schedule. Cross-zone verification would reclassify this event by requiring a second, independent
sensor to trip within a short window before treating it as high-confidence. That directly reduces
the risk of a false dispatch -- committing a costly police or guard response to an event that,
statistically, is overwhelmingly likely to be a nuisance trigger rather than a genuine intrusion.
\answerlabel{2.4.B}
\item AND-logic requires both the PIR and microwave elements to trigger within the same short
window, which sharply cuts false positives because the two technologies fail independently.
Sunlight glare fools a PIR element but produces no Doppler shift, and HVAC airflow or vibration
fools a microwave element but produces no meaningful infrared delta, so a false trigger on one
element alone is suppressed. The tradeoff is a missed true detection in the scenario where a real
intruder's movement happens to trigger only one element within the window -- for instance, someone
moving slowly enough, or in a way that produces a strong infrared delta but an unusually weak
Doppler return, so the microwave element never confirms within the required window.
\answerlabel{2.4.A}
\item Two independent changes work. First, reposition the sensor (or re-angle its detection cone) away from the
window and toward the choke point it is actually meant to protect, since a PIR aimed at or near a
window will register sunlight sweeping across the glass as a rapid infrared change. Second, reduce
sensitivity or switch to a dual-technology unit requiring AND-logic confirmation, which filters out
the single-element triggers a window view is prone to producing, without disabling detection
entirely. Placement and sensitivity are independent levers -- fixing one does not fix the other, and
a well-placed sensor tuned too aggressively will still nuisance-trigger, just on different causes.
\answerlabel{2.4.B}
\item A single sensor's false alarms are dominated by the sheer number of quiet nights, not by any
flaw a slightly better accuracy figure would fix. Even a very good detector still produces a
large absolute number of false positives across enough uneventful nights, because the rare event
being detected, not the sensor's quality, dominates the arithmetic. A second, \emph{independent}
sensor changes the structure of the problem instead of nudging one number: for a false alarm to
pass both sensors, the two unrelated failure causes (sunlight glare, HVAC airflow, and so on) must
occur simultaneously, in a way that fools both technologies at once. That is dramatically
rarer than either failure occurring alone -- exactly why cross-zone and dual-technology
verification outperform a marginal accuracy improvement on a single sensor.
\answerlabel{2.4.C}
\item Verified response directly targets the flood of false dispatches that the base-rate
arithmetic guarantees any rare-event detector will generate, by moving the confirmation step
earlier in the pipeline -- before police are sent -- rather than after. It is a structural fix, not
a stricter threshold, because it does not change any sensor's sensitivity or accuracy at all. It
inserts an independent corroborating check (a private responder or camera) between the alarm
firing and the costly response, which is the same principle behind cross-zone verification, applied
at the level of an entire city's dispatch policy rather than one alarm panel.
\answerlabel{2.4.C}
\end{answerkey}

\section*{Sources}
\begin{enumerate}[label=\lbrack\arabic*\rbrack, leftmargin=2.2em]
\item \srcitem{2.4}{1}Rana Sampson, \emph{False Burglar Alarms}, 2nd ed., Problem-Oriented Guides
for Police, Problem-Specific Guides Series No.\ 5 (U.S.\ Department of Justice, Office of
Community Oriented Policing Services, 2007); Salt Lake City figures per the same guide.
\end{enumerate}
